\documentclass{article}
\usepackage[utf8]{inputenc}
\usepackage[T1]{fontenc}
\usepackage{geometry}
\geometry{
    a4paper,
    left=2.5cm,
    right=2.5cm,
    top=2.5cm,
    bottom=2.5cm
}
\usepackage{graphicx}
\usepackage{float}
\usepackage{hyperref}
\usepackage[english]{babel}
\usepackage{amsmath}
\usepackage{amssymb}
\usepackage{array}
\usepackage{booktabs}
\usepackage{caption}
\usepackage{subcaption}
\usepackage{enumitem}
\usepackage{fancyhdr}
\usepackage{titlesec}
\usepackage{url}
\usepackage{textcomp}
\usepackage{xcolor}
\usepackage{listings}
\usepackage{multirow}
\usepackage{longtable}
\usepackage{tabularx}
\usepackage{times}

% Set up page style
\pagestyle{fancy}
\fancyhf{}
\rfoot{\thepage}
\renewcommand{\headrulewidth}{0pt}

% Title formatting
\titleformat{\chapter}[display]{\normalfont\bfseries}{\chaptertitle\thechapter}{0pt}{\Huge}
\titlespacing*{\chapter}{0pt}{50pt}{40pt}

% Document info
\title{Template Document}
\author{Academic Support Agent}
\date{\today}

\begin{document}
\maketitle

\twocolumn[
\begin{@twocolumnfalse}
\begin{tabular}{p{5cm} p{12cm}}
\multirow{2}{*}{\includegraphics[scale=0.2]{LOGO_QUIMICA_3.png}}&{{\Large\bf\titulo}}\\
\vspace{0.2cm} & \vspace{0.2cm}\\
 &{\large\nombre}\\
\vspace{0.1cm} & \vspace{0.1cm}\\ 
 & {\small\it Departamento de Qu mica. Universidad Aut noma Metropolitana Unidad Iztapalapa. San Rafael Atlixco 186. Colonia Vicentina. Iztapalapa. C.P. 09340. Ciudad de M xico, M xico.}\\
{\tiny Fecha de recepci n: No escribir nada} &\\
{\tiny Fecha de aceptaci n: No escribir nada} & {\small\resumen}\\
\end{tabular}
\end{@twocolumnfalse}
\vspace{0.5cm}
]

\section{Introducci n}
Escribe los antecedentes y la motivaci n del trabajo, sobre todo debes de escribir el estado del arte de tu tema y plantear el problema que surge de este an lisis.

\section{Ecuaciones y f rmulas}
El uso de ecuaciones en LaTex es diferente al que se usa en un procesador de texto. Sin embargo, LaTex es vers til y muy potente con el uso de s mbolos. A continuaci n se tienen algunos ejemplos:

\begin{equation}
\left(p+\frac{n^2 a}{V^2}\right)\left(V-nb\right)=nRT.\label{ec:gases_1}
\end{equation}

Otra forma de escribir una ecuaci n dentro de un texto es de la siguiente manera $p=\left(\partial E/\partial V\right)_T$. Para m s s mbolos dentro de ecuaciones, recurre al sitio \url{https://oeis.org/wiki/List_of_LaTeX_mathematical_symbols}.

La ecuaci n (\ref{ec:gases_1}) es importante para la descripci n de gases no ideales. Para una reacci n qu mica podemos escribir f rmulas qu micas \ce{H2O}, \ce{SO4^2-}, \ce{BaSO4 v}.
Y una reacci n qu mica puede ser escrita de la siguiente manera

\begin{center}
\ce{CH4 ->[O2][\Delta] CO2 + H2O}
\end{center}

En la siguiente direcci n encontrar s toda la informaci n relacionada con el paquete mhchem: \url{https://www.ctan.org/pkg/mhchem}

\section{Referencias}
Un art culo cient fico se sustenta en trabajo previo y por lo tanto el uso de referencias es muy importante. Las citas en nuestra revista se har n de acuerdo al uso de bibtex integrado en LaTex. En el archivo referencias se presentan algunos ejemplos t picos que se usa en nuestra literatura. Para citar alguna referencia que se encuentra en nuestro archivos referencia solamente hay que usar $\setminus$\texttt{cite\{Referencia\}}, por ejemplo una cita al profesor R. Sagar \cite{Sagar_1992} o la cita a un libro \cite{Szabo}. Con bibtex no tenemos que preocuparnos por la numeraci n de las referencias. Por lo tanto solamente tenemos que preocuparnos por escribir las referencias en el formato mostrado en el archivo \texttt{referencias.bib}.

\section{Figuras}
Las figuras deben ser insertadas en LaTex en formato png o eps. En la Figura \ref{fig:prueba} se presenta un ejemplo.
\begin{figure}[h]
    \centering
    \includegraphics[width=8cm]{Tlecaxitl_logo.png}
    \caption{Aqu  debemos de poner la descipci n de la figura.}
    \label{fig:prueba}
\end{figure}
En este caso se utiliza el ancho m ximo permitido en nuestra revista (8 cm).
\section{Tablas}
El reporte de resultados a trav s de tablas es muy importante para un art culo cient fico. Sugerimos los siguientes formatos de tablas. En la Tabla \ref{tab:uno} se presenta un formato donde un n mero se escribe como una ecuaci n matem tica ya que usa el s mbolo \$ para delimitar el n mero 1.61. En este caso no hay diferencia alguna si no se escribe el n mero entre este s mbolo o no. El ejemplo  reporta una tabla con tres columnas, definidas cuando se escribe el comando tabular. Las primeras dos columnas est n justificadas hacia la izquierda y la tercera se encuentra centrada. Entre las tres columnas se declara una raya vertical, todo esto se hace con el comando $\setminus$\texttt{begin\{tabular\}\{l|l|c\} (left, left y center)}.
En ocasiones es necesario hacer referencia a resultados ya publicados y que se encuentran en el archivos. En la Tabla \ref{tab:uno} se presenta un ejemplo de este tipo. 

\begin{table}[h]
\caption{Este es un ejemplo de tabla en nuestra revista.}
\label{tab:uno}
\begin{center}
\begin{tabular}{l|l|c}
\hline\hline
Nombre      &   Apellidos       &   Altura (m)\\
\hline
Enrique     &   Bonilla Cerros  &   $1.61$ \\
Carolina    &   Paredes Heredia &   1.71$^a$\\
\hline\hline
\end{tabular}
\end{center}
$^a$ Ref. \cite{1973_MEP_1}.
\end{table}

\section{Conclusiones}
Para generar documentos LaTex existe un gran n mero de ejemplos que pueden ser encontrados en muchas p ginas web o en YouTube. Lo m s importante es que este documento te debe servir para armar todo tu art culo, por favor no incluyas \texttt{package} adicionales.

\section*{Agradecimientos}
En esta secci n escribe los agradecimientos a las personas o instituciones que te permitieron desaroolar este trabajo.
\end{document}